\documentclass[11pt]{article}
\usepackage[margin=1in]{geometry}
\usepackage{amsmath,amssymb,amsthm,microtype}
\usepackage[hidelinks]{hyperref}
\newtheorem{theorem}{Theorem}
\newtheorem{lemma}[theorem]{Lemma}
\newtheorem{proposition}[theorem]{Proposition}
\newcommand{\Z}{\mathbb Z}
\newcommand{\K}{K_{\mathrm{PS}}}
\title{Parallel short trajectories in regular polygons:\\
a type convention and signed angle law}
\author{Alec Kriebel\thanks{ORCID: \url{https://orcid.org/0009-0001-9320-500X}.}}
\date{15 October 2026}
\begin{document}
\maketitle
\begin{abstract}
We give a checked, self-contained geometric proof of the signed type law
in Fuchs's Conjecture~2.6, after explicitly retaining the segment-parity
convention in his definition of type. The proof includes existence of
the required first-vertex trajectories, all polygon sizes, and boundary
directions. An unqualified type-zero shorthand is incompatible with
the unrestricted law, as direct diagonals show. Ferudun's preprint of
9~October~2026 already proves the same theorem for $n\ge4$, identifies
the same type convention issue, and uses closely related angular,
cylinder, and descent arguments. Those results are explicitly credited;
this note makes no claim of a new theorem, historical priority, or
certified independent discovery. Extensive AI assistance was used.
This is an unrefereed research note without conventional human peer review.
\end{abstract}

\section{The precise statement and its convention}
Put $N=n-2$, $\rho=\pi/n$, and $\delta=N\rho$, with $n\ge3$.
Let $P_0=O,P_1,\ldots,P_{n-1}$ be the clockwise vertices of the regular
polygon above the horizontal side $OP_{n-1}$.
A short trajectory stops at its first subsequent vertex; its open
trajectory contains no vertex. Sides are included. At its first vertex
$\alpha$ is the angle from the counterclockwise outgoing boundary
edge to the trajectory. At its last vertex the corresponding interior
angle is $\pi-\beta$, as in Fuchs's Figure~8~\cite{Fuchs}. Thus
\[
0\le\alpha\le\delta,\qquad 2\rho\le\beta\le\pi.
\]
For $q$ straight segments define the parity-selected type
$\K(\gamma)=k\in\Z/N$ by
\begin{equation}\label{type}
\begin{array}{lll}
q\text{ odd}:&s=(\alpha+\beta)/(2\rho)\in\Z,&k=n-1-s\pmod N,\\
q\text{ even}:&s=(\beta-\alpha)/(2\rho)\in\Z,&k=s-1\pmod N.
\end{array}
\end{equation}
Integrality and the geometric meaning are proved below. These formulas
retain the parity instruction immediately before Definition~2.1 of
\cite{Fuchs}. Its branch inequalities then choose representatives:
the sum case has $1\le s\le n-1$ and $0\le n-1-s\le N$.
In particular the lower side uses the representative $N\equiv0$,
before reduction, and the upper side uses $0$. In the difference case
the endpoint bounds similarly give the two difference branches.
The formulas agree with Figure~9: $OP_j$ has type $j-1$.

We explicitly adopt this completion of the printed convention.
A parity-free union of its angle branches can be ambiguous on
symmetry directions. We do not assert an authenticated author erratum
or a uniquely compelled interpretation of every printed shorthand.
Ferudun gives the same formulas and qualification in
Definition~1.1 and Remark~1.2~\cite{Ferudun}.

\begin{proposition}\label{obstruction}
For every $n\ge5$, unique types satisfying both the unrestricted
characterization $A_0\iff\beta-\alpha=2\rho$ and the unrestricted signed
law in Theorem~\ref{law} cannot exist.
\end{proposition}
\begin{proof}
The direct chord $OP_j$ has
$\alpha=(n-1-j)\rho$ and $\beta=(n+1-j)\rho$.
Every such chord therefore satisfies $\beta-\alpha=2\rho$,
although the parity-selected type is $j-1$.
For $j=n-2$, the bare characterization gives $k=0$ and
$\alpha=\rho$. The allowed choice $\ell=2$, $\varepsilon=-1$
fixes this same first-vertex ray, but the signed law would assign
$-2\ne0$ in $\Z/N$.

One can use distinct interior diagonals instead. Starting from this
chord, take $(\ell,\varepsilon)=(1,1)$ when $n$ is odd, or $(2,1)$
when $n$ is even. The resulting chords have slopes $2\rho$ and
$3\rho$, respectively, and again satisfy the bare type-zero criterion.
The required type is instead $-1$ or $-2$, nonzero modulo $N$.
Strict convexity ensures that these chords have no earlier vertex.
\end{proof}
These examples concern the unrestricted imported formulation
AMR-049-0005. Their starting angle $\rho$ is outside the odd open
sector $(0,\rho/2)$ and on the boundary of the even open sector
$(0,\rho)$ preceding the conjecture in \cite[\S2.5]{Fuchs}.
They do not refute the conjecture restricted to those generic sectors.

\begin{theorem}[Signed angle law, with the stated type]\label{law}
Suppose a short trajectory from $O$ has slope $\alpha$ and type
$\K=k$. If
\[
\alpha'=\ell\rho+\varepsilon\alpha\in[0,\delta],
\quad \ell\in\Z,\quad \varepsilon\in\{1,-1\},
\]
with $\ell$ even when $n$ is even, there is a short trajectory from $O$
with this slope, and its type is $\varepsilon k-\ell$ in $\Z/N$.
\end{theorem}
For $n\ge4$ this is also \cite[Theorem~B]{Ferudun}.
The triangle is the elementary one-type case.
The following proof uses classical double-polygon geometry
\cite{Veech,Finster}, affine centrality~\cite{BL}, and a common
labelling of outgoing and backward endpoint germs. Its reflection
pairing and descent also substantially overlap
\cite[\S\S2--3]{Ferudun}; they are not presented as new results.

\section{Endpoint labels and affine transport}
Take two distinct labelled polygons $P,Q=-P$, and glue each side to its
parallel partner by translation, reversing endpoint order.
This is the double surface $X_n$, including for even $n$.
The corner succession is $P_0,Q_1,P_2,Q_3,\ldots$.
For odd $n$ it gives one cone of angle $2N\pi$ and genus $(n-1)/2$;
for even $n$ it gives two cones of angle $N\pi$ and genus $(n-2)/2$.
These counts follow from $V-E+F$, with $E=n,F=2$.
All vertex points are retained as marked endpoints.

Number corners in a cone by $t$, beginning at $P_0$, or at $Q_0$ for
the second even cone. Give corner $t$ the angular interval
$[t\delta,(t+1)\delta]$ with coordinate $\phi$.
For odd $n$ physical direction is $\phi\bmod2\pi$.
For even $n$, cone $c\in\{0,1\}$ has physical direction
$\phi+c\pi\bmod2\pi$.
The physical lower side at $P_j$ is $-2j\rho$, and at $Q_j$ is
$\pi-2j\rho$, which checks this calibration.

Fix a real physical direction $\theta$. For either a forward germ
in direction $\theta$ or a backward germ in direction $\theta+\pi$ set
\begin{equation}\label{label}
L_\theta(\phi)=(\phi-\theta)/\pi\pmod N .
\end{equation}
For each oriented direction the $N$ germs are labelled bijectively.
For odd $n$, $N$ is odd and their labels are $2j$ or $2j+1$ modulo $N$.
For even $n$, each cone supplies $N/2$ labels, with complementary
parities at the two cones. Cone periods make \eqref{label} well defined.
For a saddle connection write $a,b$ for its initial and backward
terminal labels, with the same reference $\theta$.

Ordinary polygon development alternates translates of $P$ and $Q$.
Indeed reflection across a side equals the half-turn about its midpoint
composed with a polygon symmetry. Translation into $X_n$ therefore
identifies first-vertex trajectories with saddle connections in the
specified initial germ. The terminal copy is $P$ for odd $q$ and $Q$
for even $q$.

For a trajectory from $P_0$ take $\theta=\alpha$, so $a=0$.
At terminal corner $t$, the backward coordinate is
$t\delta+\pi-\beta$ in $P$, or $t\delta+\beta-2\rho$ in $Q$.
The physical directions give $t+s\equiv0\pmod n$ in the first case
and $t-s+1\equiv0\pmod n$ in the second, and imply the integrality of
$s$ in \eqref{type}. Substitution gives
\[
b=\begin{cases}
1+(tN-2s)/n\equiv n-1-s\pmod N,&q\text{ odd},\\
(tN+2s-2)/n\equiv s-1\pmod N,&q\text{ even}.
\end{cases}
\]
Thus $b-a=\K$; the terminal supplement and parity are essential.

\begin{lemma}\label{shift}
For $n\ge5$, every orientation-preserving affine automorphism of $X_n$
adds a common residue to all forward and backward labels in a fixed
direction. In particular it preserves $b-a$.
\end{lemma}
\begin{proof}
An angular lift $F$ of its derivative satisfies
$F(x+\pi)=F(x)+\pi$ by linear antipodality.
For odd $n$, the map on the sole cone differs from $F$ by $2\pi r$,
which gives the common shift $2r$.

For even $n$, synchronization needs justification. The half-turn
$\iota$ exchanging $P,Q$ exchanges the two cones with identical $\phi$.
It fixes exactly the $n$ glued edge midpoints. For odd $n$ it also fixes
the sole vertex point. Hence for $n\ge5$ it has $2g+2$ fixed points,
and Riemann--Hurwitz gives a sphere quotient.
Its affine conjugate is holomorphic (its derivative is $-I$) and again
has a sphere quotient. Uniqueness of the hyperelliptic involution
gives commutation with every affine automorphism, the classical
centrality statement in \cite[Lemma~2.3]{BL}.
For completeness, two distinct degree-two sphere quotients would give
a birational curve of bidegree $(2,2)$ in
$\mathbb P^1\times\mathbb P^1$, whose normalization has genus at most one.
This proves the required uniqueness in genus at least two.

If cone $0$ maps to cone $e$, its angular map is
$G(\phi)=F(\phi)-e\pi+2\pi r$ modulo $N\pi$.
Commutation with $\iota$ makes the map on cone $1$ the same $G$.
With $\theta'=F(\theta)$, every germ considered has
$\phi-\theta\in\pi\Z$, so its new label is
$L_\theta(\phi)-e+2r$ in both cones and both direction signs.
\end{proof}

\section{A finite-length reduction}

The double-polygon cylinders and shears below are classical
\cite[\S5]{Veech}; see also \cite[\S3.1]{Finster}.
The descent has the same mechanism as \cite[Lemmas~2.8--2.13]{Ferudun}.
Finite-length descent is already explicit for the pentagon in
\cite[Theorem~2.11]{DL}.
We give details for the actual double surface in both even families.

Temporarily center $P$ with vertices
$v_j=(\cos 2j\rho,\sin 2j\rho)$, and put $m=\lfloor n/2\rfloor$.
For $1\le j\le m$, let
\[
a_j=\cos 2j\rho,\quad b_j=\cos 2(j-1)\rho,\quad S_j=a_j+b_j.
\]
The vertical strip $a_j<x<b_j$ in $P$ and the reflected strip
$-b_j<x<-a_j$ in $Q$ form a cylinder. Crossing the upper edge of
$P$ translates $x$ to $x-S_j$ in $Q$; crossing the next upper edge
returns to $P$ with $x$ restored. The two vertical chord lengths
are complementary linear functions. Their sum and the transverse width are
\[
\begin{split}
\Lambda_j&=2\{\sin 2j\rho+\sin 2(j-1)\rho\}
          =4\sin((2j-1)\rho)\cos\rho,\\
H_j&=b_j-a_j=2\sin((2j-1)\rho)\sin\rho .
\end{split}
\]
Thus $\Lambda_j/H_j=\lambda:=2\cot\rho$ in every cylinder.
In flat cylinder coordinates $(x,s\bmod\Lambda_j)$, the maps
$(x,s)\mapsto(x,s+\lambda(x-a_j))$ twist once.
They fix the boundaries pointwise and have the same derivative, so glue
to one affine automorphism. This constructs the shear directly on
the double, without assuming a lift from the smaller even quotient.
Rotation by $\rho$ exchanges $P,Q$ for odd $n$; rotation by $2\rho$
preserves them for even $n$. Both respect the gluing.
Rotation of this centered picture by $-\pi/2-\rho$ identifies $v_j$
with the source vertex $P_{-j}$; the vertical model becomes direction
$-\rho$. Hence the model directions used here are source directions
$h\rho$.

\begin{lemma}\label{reduce}
For $n\ge5$, one orientation-preserving affine automorphism maps
any saddle-connection direction of $X_n$ to a direction $h\rho$.
\end{lemma}
\begin{proof}
There are finitely many saddle connections of length at most any $B$.
Indeed the complete simply connected metric cover of the unpunctured
compact Euclidean cone surface is proper and CAT(0): its cone angles
are at least $2\pi$, and the usual disk-bigon Gauss--Bonnet argument
gives uniqueness of lifted geodesics.
The finite marked endpoint set has finitely many lifts in each bounded
ball. Fixing a lift of each possible initial vertex therefore leaves
finitely many endpoint pairs at distance at most $B$, and at most one
geodesic per pair. This argument does not assume a discrete holonomy
module.

Conjugate the coordinates so that the shear has derivative
$T=\left(\begin{smallmatrix}1&\lambda\\0&1\end{smallmatrix}\right)$.
Using the actual rotations above, a saddle vector $v=(x,y)$ can be
put in a sector with $x>0$ and angle $|\eta|\le\rho/2$ for odd $n$,
or $|\eta|\le\rho$ for even $n$.
Stop if $y=0$, or at an even boundary $|\eta|=\rho$.
Otherwise an actual $T$ or $T^{-1}$ sends $v$ to
$v'=(x-\lambda|y|,y)$, and
\[
|v'|^2-|v|^2
=\lambda|y|(\lambda|y|-2x)<0
\]
The strict inequality is $x/|y|>\cot\rho$.
Rotate back into the sector and repeat. Finiteness below the initial
length forces termination. The terminal direction is the model
direction or, in the even case, one separated from it by $\rho$.
This is a finite composition of maps of the whole surface.
\end{proof}

\section{Reflection pairing and the signed law}
In a model direction $\theta=h\rho$, reflection across the polygon
symmetry axis perpendicular to $\theta$ pairs vertices by direct chords.
At a fixed vertex that direction is supporting, with no inward germ.
All other inward germs reach the reflected vertex; boundary germs
give sides and their two corner descriptions are identified by gluing.

For a chord from $P_j$ to $P_{j'}$ (or in $Q$), let
$u\in\{1,\ldots,n-1\}$ represent $j'-j$ modulo $n$.
Its initial and backward corner offsets are $(n-1-u)\rho$ and
$(u-1)\rho$. Their sum is $\delta$.
With corner numbers $t,t'$, physical alignment gives
$t+t'+1\equiv n-h\pmod n$. Thus its endpoint labels obey
\begin{equation}\label{modelsum}
a+b=\frac{(t+t'+1)N-2h}{n}\equiv-h\pmod N .
\end{equation}
This works in both polygons and for side germs.
By Lemma~\ref{reduce} and the common shift in Lemma~\ref{shift},
every direction $\theta$ containing a saddle connection has one residue
$C_\theta$ such that every connection in it satisfies
\begin{equation}\label{reflection}
b=C_\theta-a\pmod N.
\end{equation}
Every outgoing germ in that direction terminates: its image is a model
chord, whose inverse preserves the absence of intermediate vertices.
This proves the all-germ existence assertion, rather than assuming it.
The equivalent endpoint-reflection transport appears in
\cite[Lemma~3.1]{Ferudun}.

\begin{proof}[Proof of Theorem~\ref{law}]
First take $n\ge5$. The given trajectory has initial label $0$
and backward label $k$, so $C_\alpha=k$.
Suppose $\varepsilon=1$.
The isometry rotating by $-\ell\rho$ exists on $X_n$ under the stated
parity condition. Apply it to the germ at $P_0$ of slope $\alpha'$.
It now has direction $\alpha$, so terminates by the preceding argument;
pulling back proves existence at $\alpha'$.

If its new initial corner number is $t$, its preserved corner offset
is $\alpha'$ and its label is the integer
\[
a=(tN+\ell)/n,\qquad na=tN+\ell,\qquad 2a\equiv\ell\pmod N.
\]
The first integrality follows from actual physical alignment.
The last congruence uses $n=N+2$ and does not divide by two modulo $N$.
Equation~\eqref{reflection} gives $b=k-a$, hence intrinsic type
$b-a=k-\ell$. Affine invariance and the type calculation in
Section~2 identify this with the desired type.

Bisector reflection sends
$(\alpha,\beta)$ to $(\delta-\alpha,\pi+2\rho-\beta)$, preserving $q$.
It sends $s$ in \eqref{type} to $n-s$ for odd $q$, or to $2-s$
for even $q$, and so sends $k$ to $-k$.
For $\varepsilon=-1$ write
\[
\alpha'=\ell\rho-\alpha
=\delta-\{\alpha+(N-\ell)\rho\}.
\]
The bracketed angle is admissible whenever $\alpha'$ is.
Apply the positive rule with $N-\ell$, then bisector reflection.
The type becomes $-[k-(N-\ell)]=-k-\ell\pmod N$.
For even $n$, $N-\ell$ is even, as required.

For $n=3$ the sole residue is trivial. A triangular-lattice development
has a primitive first-vertex endpoint in every transformed direction,
since the allowed rotations and reflections preserve that lattice.
For $n=4$, the primitive square-lattice endpoint $(p,q)$ has type
$1$ precisely when both coordinates are odd, and type $0$ otherwise.
For positive coprime coordinates the segment count is $p+q-1$:
when both are odd, $\alpha+\beta=\pi$ gives type $1$;
when exactly one is even, the selected difference is $\pi/2$,
giving type $0$. Sides have type $0$ directly.
The interior nonidentity admissible transformation swaps coordinates;
the remaining boundary transformations exchange or fix sides.
This proves the rule and existence modulo two.
The geometric proof above already includes $n=6$; Fuchs also supplies
its elementary lattice proof in \cite[pp.~511--512]{Fuchs}.
\end{proof}

\section{Credit, chronology, and verification}
Fuchs formulated the problem and defined the types. Veech's
double surfaces and uniform cylinder shears, and classical
hyperelliptic affine centrality, are underlying inputs.
Ferudun's dated preprint~\cite{Ferudun}, Definition~1.1,
Remark~1.2 and Theorem~B, contains the same type convention,
diagonal qualification, all-admissible signed law and first-vertex
existence. Its Lemmas~2.8--2.13 and~3.1 substantially overlap the
cylinder, finite-descent, uniform angular and reflection mechanisms
used here. The full-group angular formulation above follows from
classical centrality; it is not claimed as a new discovery.
Nor is the trivial triangle case claimed as a new extension.
This note is a credited rederivation and verification of the
precisely stated result.

The fixed original candidate in
\href{https://github.com/AlecKriebel/Math/pull/191}{Math PR191}
already contains the common-germ argument.
GitHub reports PR creation on 1~October~2026, and the currently
retrieved Zenodo record for \cite{Ferudun} reports creation on
9~October~2026. Commit timestamps and present fixed-byte retrieval
do not prove when the exact proof was publicly accessible, awareness
of other work, or independent discovery. We describe these as
contemporaneous overlapping treatments, without certifying priority.
The supporting audit records source versions, hashes, read depths,
and the limits of these chronology checks.

The verification package contains exact rational and quadratic-field
first-hit tracers, the protected original candidate and checkers,
independent analytic audits, and machine-readable provenance.
The original author and original reviewer programs reproduce their
finite check counts in ordinary and optimized Python. The supplementary
type checker distributed here uses explicit failure guards.
Independent physical tracing compares 1,280 transformed pairs for
$n=3,4,5,6,8$, with 1,628 terminated trace invocations
(including repetitions). Deliberately wrong signs and terminal-angle
conventions are rejected. These computations corroborate the argument;
they do not establish universality or priority.

AI tools were used extensively for reconstruction, derivation, literature
search, drafting, computation and adversarial checking.
Separate AI audit passes are not conventional human peer review.
This preprint is unrefereed, and no conventional human peer review is
claimed at publication. No theorem in this note is advertised as
previously unknown.

\begin{thebibliography}{9}
\bibitem{Fuchs}
D.~Fuchs, \emph{Billiard Trajectories in Regular Polygons and Geodesics
on Regular Polyhedra}, Arnold Math.\ J.\ \textbf{7} (2021), 493--517.
\href{https://doi.org/10.1007/s40598-020-00170-8}
{doi:10.1007/s40598-020-00170-8}.

\bibitem{Ferudun}
A.~Ferudun, \emph{On Seven Conjectures of Fuchs on Billiard Trajectories
in Regular Polygons and Geodesics on the Regular Dodecahedron},
unrefereed preprint, version~1.0, 9~October~2026.
\href{https://doi.org/10.5281/zenodo.23267641}
{doi:10.5281/zenodo.23267641}.

\bibitem{Veech}
W.~A.~Veech, \emph{Teichm\"uller curves in moduli space, Eisenstein
series and an application to triangular billiards},
Invent.\ Math.\ \textbf{97} (1989), 553--583.
\href{https://doi.org/10.1007/BF01388890}{doi:10.1007/BF01388890}.

\bibitem{Finster}
M.~Finster, \emph{A series of coverings of the regular $n$-gon},
\href{https://arxiv.org/abs/1005.4588v3}{arXiv:1005.4588v3} (2011),
\S\S3.1--3.2.

\bibitem{BL}
C.~Boissy and E.~Lanneau, \emph{Pseudo-Anosov homeomorphisms on
translation surfaces in hyperelliptic components have large entropy},
\href{https://arxiv.org/abs/1005.4148v2}{arXiv:1005.4148v2} (2010),
Lemma~2.3. The earlier author draft of 21~July~2010 numbers the
same centrality statement Lemma~2.2.
\bibitem{DL}
D.~Davis and S.~Leli\`evre, \emph{Periodic paths on the pentagon,
double pentagon and golden L},
\href{https://arxiv.org/abs/1810.11310v1}{arXiv:1810.11310v1} (2018),
Theorem~2.11.
\end{thebibliography}
\end{document}
